
\subsection{Dogmatic Paragraphs on the Church}

Georg Sverdrup
Manuscript 1907

§ 1. To treat the doctrine of the Church after the doctrine of the means of grace, so that the Church should be “the community of people among whom the means of grace have had their intended effect” (Tønning), does not answer to the Christian understanding of the Church’s coming into being, nor either to the individual Christian’s experience of how he himself has been incorporated into the Church.

§ 2. The Church came into being through a divine act, which is part of God’s revelation among men, and belongs to the eternal counsel of salvation which promises the race’s restoration from its fall. This divine act consists in the sending of the Holy Spirit into the hearts of Jesus’ disciples on the first (Jewish) Pentecost after Christ’s ascension. The signs which accompanied the Spirit’s sending are such as reveal what the Spirit’s working is like in character, and what it aims at. The storm signifies the mighty power, the tongues of fire the heavenly light, the foreign tongues the spread among all peoples.

§ 3. The divine benefaction which the Spirit’s sending brings to Jesus’ believing disciples is, according to Jesus’ promise in the Gospel of John chapters 14-17, Luke 24:49, and Acts 1:4-8, and according to Christian experience of the Spirit’s working: new, divine life, and new, spiritual light joined with the power for action and urge to witness that belong thereto. Both the divine life, which according to its essence is love, and the spiritual light, which is truth, proceed from the exalted and glorified Jesus Christ, with whom the Spirit unites the believing disciples, so that they become His body, and He their Head.

§ 4. The Church is thus, for Christian understanding, a new creation, a new race, which again bears the image of God; yet not a new bodily human race alongside the old, but a renewal of the old race through the forgiveness of sins and the life of the Holy Spirit in the hearts, in such a way that this community lives the renewed human life amid the outward conditions of the fallen race, surrounded by sin in the flesh and in the world, and with temporal death as the regular, though not unavoidable, issue from the present world.

§ 5. The Church, which consists of those who have received the Holy Spirit, takes on its outward form, manifests itself, and shows itself effectual and capable of inward and outward growth through the use of the gifts of grace and the means of grace. All have received the same Spirit, but not therefore precisely the same gifts; therefore there arise several different branches of activity in the Church through such persons as God has given gifts, and the Church has received by its recognition of the gifts. Of all activity in the Church, the use of the means of grace, the Word and the Sacraments, is the most significant and, under all historical conditions, necessary for the Church’s preservation and growth; and for that reason the gifts of grace which fit men for the administration of the means of grace, and most particularly for the proclamation of God’s Word, are those which the Church esteems most highly and cultivates with particular care.

§ 6. The Spirit’s working for the Church’s inward and outward growth thus takes place through the gifts of grace and the means of grace, and since order and coherence are needed in the Church’s work within its own midst and in and for the surrounding world, there follows from this of necessity the Church’s organization, in order that the work may become as effective as possible, and all members of the Church may find their place in it.

§ 7. The relation of the means of grace to the Church consists in this, that by the Word faith in Jesus Christ is created and preserved; by Baptism the believing are received into living fellowship with Jesus Christ and thereby into His body, the Church; and by the Lord’s Supper this living fellowship between the Lord and the Church is sustained and preserved.

§ 8. The relation of the Church to the means of grace consists in this, that the Church administers them through men whom it elects for that purpose; in this its election the Church seeks to find the men whom the Lord Himself has furnished with the gifts which are especially required for the proclamation of the Word and the Church’s spiritual guidance. The “inner call” to this work thus consists essentially in the gift of grace from God, whereas the “outer call” consists in the Church’s election to labor in the definite field. Since no difference of rank takes place among the Church’s members, the Church’s election or ordination cannot produce any such difference either, since the matter concerns only the distribution of work among the persons fitted for it. The less therefore Scripture permits the priests to rule in the Church, the more it holds true that all in the Church owe the obedience of faith to the Lord’s Word.

§ 9. From these outward conditions of the Church it follows that from the one first Church many congregations spring forth, which nevertheless all together constitute only the one Christian community on earth, which began with the sending of the Holy Spirit, which community therefore also in the New Testament is called by the same name as the individual congregation, namely ecclesia, whereas we frequently give it the name Church (from kyriaké, an adjective with ecclesia understood, meaning “belonging to the Lord”). Yet this new name must not lead us to think that thereby there should be designated a community which stood above the congregation or had an essence different from it.

§ 10. This Christian community on earth is, despite the division into many congregations, despite the various confessions, in part mutually contradictory, despite the many different ecclesiastical organizations, despite the theological controversies, still only one, in that the unity consists (invisibly) in the one Spirit and the one Lord with all the spiritual qualities that follow therefrom, and (visibly) in the one and the same Word of God, the one Baptism, the one Lord’s Supper, the common ecumenical confessions, and a multiplicity of less significant marks. The most important division of which the history of the Church reports is that which took place at the Reformation, in that the opposition to the Church’s secularization, which was present more or less strongly in the Church history of antiquity and the Middle Ages, became so strong that it brought about a breach which can be healed only slowly. But neither does this abolish the Church’s unity, because the Reformation has precisely its roots in the Christian foundation which was found within the Roman Church itself, though buried beneath the rubble of many human ordinances.

§ 11. The one Christian community on earth is holy on account of the indwelling and working of the Holy Spirit therein, whereby the believing disciples of Jesus are sanctified in an ever more inward union with the Lord Jesus Christ, while those who have fallen away from the faith, through the working of the Word and the Spirit upon their resistant hearts, ripen unto the final and unavoidable judgment. Holiness is therefore not abolished by the participation of hypocrites in the outward Church.

§ 12. In the Apostles’ Creed the Christians also confess faith in the Church’s catholicity, which consists in this, that the Holy Spirit works for the salvation and eternal blessedness of all men, and that this missionary activity does not cease until all who will let themselves be saved have been saved, and thus the work of the Son and the Spirit on earth has reached its perfect completion, in that the fallen human race has been restored from its fall, and those who, despite the Gospel’s offer, have resisted the Spirit and grace, are excluded from the kingdom of God’s eternal glory and from the redeemed human race.

§ 13. The Church’s apostolicity, or the Christian conviction that the Church still is the same as that which in the New Testament period of revelation was built upon the foundation of the apostles and prophets, rests upon the truth that the Church’s coming into being and all its existence and history are wrought and borne by the Holy Spirit; and it is served by this, that the Church preserves agreement with the holy New Testament Scripture, which is the complete and infallible witness concerning the Church’s essence and reality from the beginning.

§ 14. Since the Church has come into being through the Spirit’s sending into the hearts of Jesus’ earthly disciples, and these have Jesus’ command to make all peoples His disciples, it becomes unavoidable that the question should be raised and decided already in the apostolic age whether the Church is only a form of the Jewish religious community modified by faith in the Messiah who has come, and whether therefore the Church is bound by the Law that was given to the Jews, so that the Church can be known only by the observance of Jewish customs and ceremonies. And through the struggle which the deciding of this question costs, the great result is won for Christian faith, that the Church is a new community dependent only on faith in Jesus Christ and the new creation, which consists in the Spirit’s imparting of the new life, so that none of the signs and symbols, customs and ceremonies of the Jewish Law are binding upon the Church and of course neither necessary for salvation. This mighty victory over all demands for the observance of the Law has the abiding significance for the Church in all ages, that it is impossible to bind it by such regulations, laws, customs, ceremonies, and the like, as are commandments of men, and make them, or the observance of them, marks of the one true Church.

§ 15. The question whether the Kingdom of God, or the Kingdom of Heaven, is the same as the Church, must first be answered to this effect, that the Church can only be the Kingdom of God on earth, and second, that even if we speak only of the Kingdom of God here on earth or in the present world, yet the concept of the Kingdom of God is broader than the concept of the Church, in that to it belongs not only the Spirit’s life in the hearts, which is the Church’s inward essence, but also the preparing and guiding divine government and governance of the human race, and doubtless also the manifold effects of divine revelation in the Gospel outside of and alongside the Church in the life of the human spirit, in art, science, literature, the constitution of the state, etc.; just as also the exceedingly great changes which take place in the world, and especially in the existence of the human race, at Christ’s return belong under the concept of the Kingdom of God. The relation between the Church and the Kingdom of God on earth must therefore be determined to this effect, that the Church is the central current of life in the Kingdom of God on earth, alongside of which divine operations of grace and power are perceived, which indeed belong to the Kingdom of God, but yet do not serve unto salvation for those who are touched by them, unless they become means to draw them to the Son and His body, which is the Church. The Church’s course is concluded and brought to completion at Christ’s return and the final judgment, when the Kingdom of glory with its rest and blessedness replaces the Church’s labor, struggle, and tribulation.

Can the Lutheran Free Church Possess Property?

At the Free Church’s annual meeting in 1897 the question whether the Free Church can possess property was, according to the report in “Foldebladet” for June 16, 1897, answered as follows:

(Prof. Sverdrup said) “the Free Church could own nothing and yet possess all things in a spiritual sense. The Free Church can no more own the bookstore than it can own Augsburg Seminary or the Deaconess Home.”
